\documentclass[10pt,twocolumn]{article}
\usepackage[T1]{fontenc}
\usepackage[utf8]{inputenc}
\usepackage{lmodern}
\usepackage[margin=1.8cm,top=2.4cm,bottom=2cm,columnsep=0.6cm]{geometry}
\usepackage{fancyhdr}
\usepackage{titlesec}
\usepackage{booktabs}
\usepackage{amsmath,amssymb,amsfonts}
\usepackage{microtype}
\usepackage{xcolor}
\usepackage{mdframed}
\usepackage{parskip}
\usepackage{enumitem}
\usepackage{hyperref}

\definecolor{rulered}{RGB}{180,30,30}
\definecolor{boxgray}{RGB}{245,245,245}
\definecolor{keywordgray}{RGB}{100,100,100}

\pagestyle{fancy}
\fancyhf{}
\fancyhead[L]{\textsc{\small Claude Mag}}
\fancyhead[C]{\small\textit{Machine Learning Research Digest}}
\fancyhead[R]{\small 2026-08-03}
\fancyfoot[C]{\small\thepage}
\renewcommand{\headrulewidth}{0.8pt}

\titleformat{\section}{\normalsize\bfseries\color{rulered}}{}{0pt}{}%
  [\vspace{-4pt}\color{rulered}\rule{\columnwidth}{0.4pt}]
\titlespacing{\section}{0pt}{8pt}{4pt}

\hypersetup{colorlinks=true,urlcolor=rulered,linkcolor=black}

\begin{document}
\setlength{\parindent}{0pt}
\setlength{\parskip}{4pt}

{\small\color{keywordgray}\textbf{Keywords:}
  agentic reinforcement learning \textbullet{}
  skill evolution \textbullet{}
  cross-task transfer \textbullet{}
  credit assignment}\par\vspace{4pt}

{\LARGE\bfseries\raggedright
  SkillRise: Agentic Reinforcement Learning for
  Cross-Task Skill Evolution}\par\vspace{4pt}

{\large\itshape\raggedright
  A Single Policy Alternates Between Solving Tasks and Curating a
  Growing Skill Document---Rewarded for Downstream Impact, Not
  Document Plausibility}\par\vspace{6pt}

{\small\textbf{By} Zhiyuan Yao, Yuxin Chen, Zhengxi Lu, Zishan Xu,
  Yueqing Sun, Yifu Guo, Yuquan Lu, Zhengzhou Cai, Kangning Zhang,
  Zhuowen Han, Zihan Wang, Ziang Ye, Qi Gu, Xunliang Cai,
  Weiwen Liu, Yongliang Shen
  (Zhejiang University; NUS; SJTU; Meituan)}\par\vspace{2pt}

{\small\color{keywordgray}arXiv:2607.26784 \textbullet{} July 29, 2026}
\par\vspace{2pt}\noindent{\color{rulered}\rule{\columnwidth}{0.6pt}}\vspace{6pt}

Standard agentic RL treats each task as an independent episode, discarding
transferable workflows after every interaction. \textbf{SkillRise} trains
a single policy to simultaneously solve tasks and maintain an evolving
natural-language skill document. A key innovation is \emph{decoupled credit
assignment}: skill curation is rewarded based on its downstream impact on
future tasks, not on document quality alone. On ALFWorld, WebShop, and
ScienceWorld, SkillRise achieves the best Pass@1 among all compared methods.

\begin{mdframed}[backgroundcolor=boxgray,linecolor=rulered,linewidth=1pt,
    innerleftmargin=6pt,innerrightmargin=6pt,
    innertopmargin=6pt,innerbottommargin=6pt]
\textbf{\small AT A GLANCE}\par\vspace{4pt}
\begin{description}[leftmargin=0pt,labelsep=4pt,itemsep=2pt,parsep=0pt,
    font=\small\bfseries]
  \item[Problem:] Agentic RL discards generalizable patterns after each
    task, preventing efficient transfer to future related tasks.
  \item[Why it matters:] Real-world agents face streams of related tasks;
    the ability to curate and reuse skills is essential for continual improvement.
  \item[Method:] Single policy that solves tasks and curates a skill
    document; decoupled credit assignment rewards curation by downstream
    task improvement.
  \item[Result:] Best Pass@1 on ALFWorld (78.9\%), WebShop (66.3\%),
    ScienceWorld (50.2\%).
  \item[Evidence:] 11--12 pp. gain over standard RLVR; shared solve-curate
    policy outperforms separate-policy baselines.
\end{description}
\end{mdframed}

\section{The Big Picture}

Agentic reinforcement learning for LLMs (RLVR) has demonstrated that
language agents can learn to navigate web interfaces, plan household tasks,
and execute scientific workflows purely from outcome-based rewards. But these
agents are trained and tested on episodic tasks: each episode is independent,
and the knowledge gained in one episode is only preserved insofar as
it is absorbed into the model's parameters through gradient updates.

This means every new task starts from scratch. A household agent that learned
to find items in a drawer on task 1 has no mechanism to record that workflow
and apply it directly on task 2. SkillRise proposes a simple but effective
alternative: let the agent write its own workflow documentation, and reward
the documentation by how much it helps future tasks.

\section{Why Existing Approaches Fall Short}

Prior work on skill-augmented agentic RL (Voyager, SKILL1, SkillRL,
SkillFlow) decomposes skill evolution into a pipeline:
\begin{enumerate}[itemsep=1pt,parsep=0pt]
  \item \textit{Extract:} A separate skill extraction module reads
    completed trajectories and identifies transferable patterns.
  \item \textit{Store:} Extracted skills are written to a retrieval database.
  \item \textit{Retrieve:} At test time, a retriever fetches relevant
    skills as context.
\end{enumerate}

This pipeline has two weaknesses. First, it requires separate optimization
objectives for each component: extraction, retrieval, and execution are
trained separately with potentially conflicting objectives. Second, the
skill document is static after extraction and cannot be revised in light of
future task failures---a pattern extracted from an easy instance may
mislead on a harder variant.

SkillRise collapses all three stages into a single reinforcement learning
objective applied to a single policy.

\section{Core Mechanism}

SkillRise operates over a sequence of $N$ tasks $T_1, T_2, \ldots, T_N$
organized from easy to hard. For each task $T_t$, the agent is given the
current skill document $S_{t-1}$ and takes two sequential actions:

\textbf{Solve:} Generate a trajectory $a_{1:H}$ to complete $T_t$ given
context $(T_t, S_{t-1})$.

\textbf{Curate:} Given $(T_t, a_{1:H}, S_{t-1})$, generate an updated
skill document $S_t$ that adds, revises, or deletes workflow entries
in light of what was learned on $T_t$.

Both actions are produced by the same policy $\pi_\theta$; no separate
extraction or retrieval module exists.

\section{Technical Formulation}

Let $r_{\mathrm{solve}}(T_t, a_{1:H})$ be the task completion reward and
$r_{\mathrm{curate}}(S_t, T_{t+1:t+k})$ be the discounted impact of the
updated skill document on the next $k$ tasks:

\[
  r_{\mathrm{curate}}(S_t, T_{t+1:t+k}) =
    \frac{1}{k}\sum_{j=1}^{k} \gamma^j \cdot
    \Delta r_{\mathrm{solve}}(T_{t+j},\; S_t \text{ vs.\ } S_{t-1})
\]

The SkillRise objective is:
\[
  \mathcal{J}(\theta) = \mathbb{E}\!\left[
    \sum_{t=1}^{N} r_{\mathrm{solve}}(T_t, a_{1:H}^{(t)})
    + \gamma \cdot r_{\mathrm{curate}}(S_t, T_{t+1:t+k})
  \right]
\]

This is the critical distinction from prior work: curation is rewarded by
downstream impact ($\Delta r_{\mathrm{solve}}$), not by surface-level quality.
A skill is only reinforced when it demonstrably helps solve future tasks.

\section{Learning or Inference Procedure}

Training uses Group Relative Policy Optimization (GRPO): multiple curation
outputs are sampled for each post-task state, and the policy is updated
based on relative advantage between outputs that have better vs.\ worse
downstream impact.

The training curriculum organizes tasks by difficulty within each domain:
easier instances first. This ensures that skill documents accumulate useful
content early in training before harder tasks demand more nuanced workflows.
Skill documents are constrained to 2,000 tokens to fit within the context
window alongside task instructions.

\section{What the Guarantee Says}

Decoupled credit assignment provides a monotonicity property: if the
curation policy is improved, the downstream solve reward is non-decreasing
in expectation. This follows from the definition of $r_{\mathrm{curate}}$:
if $S_t$ improves performance on future tasks, the curation reward is
positive, incentivizing the policy to generate more useful updates.

For a fixed backbone, the SkillRise objective upper-bounds the cumulative
task completion rate over the $N$-task stream.

\section{Experimental Findings}

\textbf{ALFWorld (household task completion, Pass@1):}

\begin{center}\small
\begin{tabular}{lc}
\toprule
Method & Pass@1 \\
\midrule
Base LLM & 42.3\% \\
Standard RLVR & 67.8\% \\
SKILL1 & 71.4\% \\
SkillRise & \textbf{78.9\%} \\
\bottomrule
\end{tabular}
\end{center}

\textbf{WebShop (e-commerce navigation, Task Success):}

\begin{center}\small
\begin{tabular}{lc}
\toprule
Method & Success \\
\midrule
Standard RLVR & 56.2\% \\
SkillFlow & 61.4\% \\
SkillRise & \textbf{66.3\%} \\
\bottomrule
\end{tabular}
\end{center}

\textbf{ScienceWorld (scientific experiment tasks, Pass@1):}

\begin{center}\small
\begin{tabular}{lc}
\toprule
Method & Pass@1 \\
\midrule
Standard RLVR & 38.1\% \\
SKILLC & 43.7\% \\
SkillRise & \textbf{50.2\%} \\
\bottomrule
\end{tabular}
\end{center}

Gains over standard RLVR are 11.1\,pp (ALFWorld), 10.1\,pp (WebShop),
and 12.1\,pp (ScienceWorld).

\section{Ablations and Interpretation}

\textbf{Decoupled vs.\ immediate curation reward:} Replacing
$r_{\mathrm{curate}}$ with an immediate quality score (document fluency
or coverage) reduces ALFWorld Pass@1 by 6.6\,pp, confirming that
downstream impact is the essential signal.

\textbf{Shared vs.\ separate policy:} Training a separate curation policy
(independent weights) reduces ALFWorld Pass@1 by 4.1\,pp, suggesting
productive weight-sharing: solving improves curation, and curation
improves solving through shared world model knowledge.

\textbf{Curriculum ordering:} Random task ordering reduces ALFWorld
Pass@1 by 6.3\,pp, confirming that the easy-to-hard curriculum is
essential for early skill accumulation before difficult tasks arrive.

\textbf{Document length:} Allowing up to 4,000 tokens (vs.\ 2,000) gains
1.2\,pp on ScienceWorld, suggesting the constraint is occasionally binding
for complex scientific workflows.

\section*{Reference}

Zhiyuan Yao, Yuxin Chen, Zhengxi Lu, et al.
\textbf{SkillRise: Agentic Reinforcement Learning for Cross-Task Skill
Evolution}.
arXiv:2607.26784, July 29, 2026.
\url{https://arxiv.org/abs/2607.26784}

\end{document}
